\chapter{INTRODUÇÃO}

Futebol é o esporte mais popular do mundo \cite{Almulla2023}. Nas últimas décadas, sua evolução tem sido notável, acompanhada por um crescimento expressivo 
nos lucros gerados pelas atividades esportivas \cite{Rodrigues2022}. No entanto, o futebol é um esporte altamente imprevisível. 
Um exemplo notável dessa imprevisibilidade foi a surpreendente conquista da English Premier League (EPL) pelo Leicester City em 2016, um feito que surpreendeu 
especialistas e torcedores em todo o mundo \cite{Baboota2019}.

Dada a relevância econômica e o impacto do futebol como forma de entretenimento, pesquisadores da academia e profissionais da indústria têm se dedicado ao 
desafio de prever seus resultados \cite{Baboota2019}. O crescimento do mercado de apostas reflete essa busca, impulsionando também os lucros da indústria de 
prognósticos esportivos \cite{Rodrigues2022}. Para se ter uma noção da magnitude econômica do futebol, entre 1º de junho e 1º de setembro de 2023, as transferências 
de jogadores movimentaram R\$ 36,6 bilhões, conforme relatório divulgado pela Federação Internacional de Futebol (FIFA) \cite{Globo2023}. No Brasil, em 2025, 
a janela de transferências da Série A brasileira bateu a marca de R\$ 1 bilhão em movimentações financeiras \cite{Globo2025}.

No campo da Ciência da Computação, muitos estudos voltam-se para a análise de dados e estatística, explorando métodos para modelar padrões e fazer previsões. 
Um exemplo relevante é o estudo de Joseph, Fenton e Neil (2006), que investigou o uso de redes Bayesianas e outras técnicas de Aprendizado de Máquinas (AM) para 
prever resultados do Tottenham Hotspur entre 1995 e 1997. Outro caso notável de aplicação de dados no futebol envolve o \textit{Brighton \& Hove Albion}, clube inglês que, 
sob a gestão de Tony Bloom, implementou uma estratégia fortemente orientada por análises estatísticas para recrutamento e tomada de decisões. Essa filosofia refletiu-se
diretamente na política de transferências do clube, com a identificação sistemática de atletas subvalorizados no mercado. Entre os exemplos mais emblemáticos está
Moisés Caicedo, contratado junto ao Independiente Del Valle em 2021 por cerca de \pounds4,5 milhões e negociado com o Chelsea em 2023 por valores próximos a 
\EUR{115 milhões}, tornando-se uma das maiores vendas da história da EPL \cite{BBC2023Caicedo, CNN2023Caicedo}. O mesmo caso de altos lucros em revendas se aplica 
a outros jogadores, como Alexis Mac Allister e Marc Cucurella, abordagem que permitiu que o \textit{Brighton} se consolidasse como um dos clubes 
mais eficientes do mercado de transferências da EPL, reinvestindo os recursos obtidos e mantendo a competitividade sustentável na elite do futebol inglês \cite{ESPN2023}.

Na literatura, diversos estudos propõem a utilização de algoritmos de AM para prever resultados, com o objetivo de otimizar ganhos financeiros 
no contexto de apostas. No entanto, poucos trabalhos focam na construção de elencos com o objetivo de aprimorar a seleção de jogadores \cite{Abidin2021}.
Observa-se, portanto, uma lacuna na literatura no que diz respeito ao uso de AM como ferramenta de apoio à construção e otimização de elencos esportivos.
A escolha de um jogador para compor um elenco envolve a análise de múltiplos parâmetros, incluindo características técnicas, físicas e mentais \cite{Abidin2021}. 
Conforme discutido anteriormente, a natureza não determinística inerente ao futebol torna essa tarefa ainda mais desafiadora. Mesmo uma contratação que atenda a todos os 
critérios desejados pode, ainda assim, não atingir o desempenho esperado.

Segundo Ian Graham (2024), uma transferência pode ser considerada um sucesso se o jogador iniciar 50\% das partidas nas primeiras duas temporadas ("The 50\% rule"). 
Considerando todas as transferências de mais de € 10 milhões que ocorreram na EPL entre os anos de 1992 e 2021, perto da metade (46\%) não atingiram o critério 
estabelecido e, portanto, podem ser consideradas como falhas \cite{Graham2024}. Alguns dos principais motivos para essas falhas são:

\begin{enumerate}
    \item O jogador atual é melhor que o novo jogador;
    \item O jogador não é tão bom quanto se pensava inicialmente;
    \item O jogador não se encaixa no estilo do time;
    \item O jogador joga fora de posição;
    \item O técnico não aprova o jogador;
    \item O jogador tem problemas de condicionamento físico ou pessoais.
\end{enumerate}

Nos últimos anos, o AM tem revolucionado a maneira como os dados esportivos são analisados, permitindo que padrões antes imperceptíveis sejam 
descobertos a partir de grandes volumes de informação. Técnicas avançadas vêm sendo aplicadas para prever desde a probabilidade de vitória em uma partida até o 
impacto de decisões táticas ao longo de um campeonato. Neste estudo, é explorado o uso de modelos preditivos com o intuito de converter estatísticas históricas 
de jogadores em insights estratégicos, auxiliando clubes nas tomadas de decisões de mercado, principalmente nas aquisições de jogadores para complementar o elenco.
Dessa forma, este trabalho propõe uma abordagem orientada por dados para apoiar decisões de mercado no futebol profissional, utilizando modelos preditivos 
capazes de estimar o desempenho esperado de jogadores em contextos específicos de equipe e temporada.